\documentclass[11pt,a4paper]{article}
\usepackage{shared/luxcover}
\usepackage[utf8]{inputenc}
\usepackage[T1]{fontenc}
\usepackage{amsmath,amssymb,amsfonts,amsthm}
\usepackage{graphicx}
\usepackage{hyperref}
\usepackage{booktabs}
\usepackage{algorithm}
\usepackage{algpseudocode}
\usepackage{enumitem}
\usepackage{xcolor}
\usepackage{listings}

\hypersetup{colorlinks=true,linkcolor=black,citecolor=blue,urlcolor=blue}

\newtheorem{theorem}{Theorem}[section]
\newtheorem{lemma}[theorem]{Lemma}
\newtheorem{definition}{Definition}[section]
\newtheorem{remark}{Remark}[section]

\lstset{
  basicstyle=\ttfamily\small,
  breaklines=true,
  frame=single,
  numbers=left,
  numberstyle=\tiny,
  language=Go
}

\title{LP-Anchor: On-Chain Checkpoint Anchoring\\for CRDT State Durability}

\author{
Zach Kelling\thanks{Corresponding author: zach@lux.network}\\
\textit{Lux Industries, Inc.}\\
\texttt{research@lux.network}
}

\date{April 2026}

\begin{document}

\luxcoverpage

\begin{abstract}
We describe a mechanism for anchoring CRDT document state to a Lux chain
so that clients hold cryptographic proof that a given Merkle root existed at
a given logical height. The anchor precompile accepts a 72-byte payload
(32-byte app ID, 8-byte height, 32-byte root) and stores it in EVM state
at approximately 25,400 gas on C-Chain or near-zero cost on Q-Chain via
Quasar DAG parallelism. Users verify anchors by querying the chain and
recomputing their local Merkle root. A rollup-style batching mode amortizes
cost across 1,000 anchors per on-chain transaction. The design is
implemented as a stateful EVM precompile at address \texttt{0x0700\ldots10}
in the Lux precompile registry (Privacy/Encryption page, item 0x10).
\end{abstract}

\tableofcontents

% ============================================================================
\section{Problem Statement}
\label{sec:problem}

Base (and any CRDT-based application) replicates state across clients via
conflict-free merge operations. The server is the coordination point: it
receives operations, merges them, and fans out updates. This creates a
trust asymmetry:

\begin{enumerate}[nosep]
  \item The server can \textbf{withhold operations}: silently drop writes
        from a subset of clients.
  \item The server can \textbf{rewrite history}: present different state
        vectors to different clients without detection.
  \item The server can \textbf{lie about time}: claim a state snapshot
        existed at time $T$ when it did not.
\end{enumerate}

\noindent
The user needs a \emph{tamper-evident durability backstop}: cryptographic
proof that the CRDT state at logical height $h$ had Merkle root $r$,
anchored to a public chain whose finality is independent of the Base server.

\begin{definition}[Anchor]
An anchor is a tuple $(\mathit{appID}, h, r, t, \sigma)$ where
$\mathit{appID} \in \{0,1\}^{256}$ identifies the application,
$h \in \mathbb{N}$ is the logical height (monotonically increasing),
$r \in \{0,1\}^{256}$ is the Merkle root of the CRDT state at height $h$,
$t$ is the chain timestamp at inclusion, and $\sigma$ is the submitter's
ECDSA signature (implicit in the EVM transaction).
\end{definition}

% ============================================================================
\section{Architecture}
\label{sec:architecture}

\subsection{Anchor Lifecycle}

\begin{enumerate}[nosep]
  \item \textbf{Compute}: Application computes SHA-256 Merkle root of
        serialized CRDT \texttt{DocumentSnapshot} every $N$ operations or
        $T$ seconds (configurable; default $N=1000$, $T=300$s).
  \item \textbf{Submit}: Application calls
        \texttt{anchor.Submit(appID, height, root)} which encodes a
        72-byte payload and sends an EVM transaction to the anchor
        precompile.
  \item \textbf{Finalize}: Transaction is included in a block and
        finalized by Quasar certificate (Q-Chain) or Snow consensus
        (C-Chain). The anchor is now immutable.
  \item \textbf{Verify}: Any party calls \texttt{anchor.Get(appID, height)}
        via \texttt{eth\_call}, recomputes the local Merkle root from their
        CRDT state, and compares.
\end{enumerate}

\subsection{Precompile Design}

The anchor precompile is a stateful EVM precompile registered at address
\texttt{0x0700\ldots10} (Privacy/Encryption page, C-Chain slot, item 0x10
in the LP addressing scheme). It exposes four functions via ABI selector:

\begin{center}
\begin{tabular}{llr}
\toprule
\textbf{Function} & \textbf{Selector} & \textbf{Gas} \\
\midrule
\texttt{submit(bytes32,uint64,bytes32)} & \texttt{0x3c21e8a1} & 25,400 \\
\texttt{get(bytes32,uint64)} & \texttt{0x8eaa6ac0} & 2,600 \\
\texttt{getLatest(bytes32)} & \texttt{0xba41d847} & 2,600 \\
\texttt{list(bytes32,uint64,uint64)} & \texttt{0x70fa9a4d} & 2,600 + 200/entry \\
\bottomrule
\end{tabular}
\end{center}

\subsubsection{Storage Layout}

Each anchor occupies two EVM storage slots:

\begin{itemize}[nosep]
  \item Slot $\texttt{keccak256}(\mathit{appID} \| h)$: the 32-byte root.
  \item Slot $\texttt{keccak256}(\mathit{appID} \| \texttt{``latest''})$:
        the highest submitted height for this appID (8 bytes, right-padded).
\end{itemize}

\noindent
This gives $O(1)$ lookup per anchor and $O(1)$ latest-height query. The
\texttt{list} function iterates heights $[h_{\text{from}}, h_{\text{to}})$,
issuing one \texttt{SLOAD} per height.

\subsubsection{Gas Cost Derivation}

\begin{align*}
\text{submit gas} &= G_{\text{base}} + 2 \times G_{\text{sstore\_new}} + G_{\text{keccak}} \\
                  &= 400 + 2 \times 12{,}000 + 600 + 400 \\
                  &= 25{,}400
\end{align*}

\noindent
where $G_{\text{sstore\_new}} = 12{,}000$ for writing a new storage slot
(after EIP-2929 warm access), $G_{\text{keccak}} = 600$ for the two hash
computations, and $G_{\text{base}} = 400$ for ABI decoding and validation.
A \texttt{get} call costs 2,600 (one warm \texttt{SLOAD}).

\subsection{Height Monotonicity Invariant}

The precompile enforces $h > h_{\text{latest}}$ for a given appID. Attempts
to submit a height $\leq$ the current latest revert with
\texttt{AnchorHeightNotMonotonic}. This prevents retroactive anchor
forgery: once height 100 is anchored, no one can insert a different root
at height 100 or below.

% ============================================================================
\section{Why Q-Chain}
\label{sec:qchain}

The Q-Chain (Quantum Chain) is the preferred anchor target for three reasons.

\paragraph{PQ-safe finality.}
Q-Chain blocks carry Quasar certificates: triple-signed (BLS + ML-DSA +
Ringtail) finality proofs. An anchor finalized by a QuasarCert inherits
post-quantum security. C-Chain finality relies on Snow consensus
(probabilistic, classical signatures only).

\paragraph{DAG parallelism.}
Q-Chain processes transactions as a DAG, not a sequential block. Anchor
submissions from independent appIDs are conflict-free and process in
parallel. On C-Chain, all transactions compete for sequential block space.

\paragraph{Cost.}
Q-Chain transaction fees are validator-set-determined and currently near
zero for data-only transactions (no EVM execution). A C-Chain anchor at
25,400 gas $\times$ 25 nLUX/gas = 635 nLUX ($\approx$ \$0.001 at
\$1.50/LUX). Q-Chain is 10--100$\times$ cheaper for equivalent data
anchoring because it avoids EVM metering overhead.

\paragraph{C-Chain alternative.}
If the application already pays C-Chain gas (e.g., a DeFi protocol on
C-Chain), using the same chain for anchoring avoids cross-chain
complexity. The precompile is registered on both chains.

% ============================================================================
\section{Anchor Transaction Format}
\label{sec:format}

The on-chain payload is 72 bytes:

\begin{center}
\begin{tabular}{lrl}
\toprule
\textbf{Field} & \textbf{Bytes} & \textbf{Encoding} \\
\midrule
appID  & 32 & \texttt{bytes32} (SHA-256 of app name) \\
height & 8  & \texttt{uint64}, big-endian \\
root   & 32 & \texttt{bytes32} (Merkle root) \\
\bottomrule
\end{tabular}
\end{center}

\noindent
Total: 72 bytes of calldata + 4-byte function selector = 76 bytes.
The submitter identity is \texttt{msg.sender} (implicit ECDSA signature
in the EVM transaction envelope). No additional signature field is needed.

% ============================================================================
\section{Verification Path}
\label{sec:verification}

A user who suspects the server is lying:

\begin{enumerate}[nosep]
  \item Queries \texttt{getLatest(appID)} to learn the highest anchored
        height $h_{\max}$.
  \item Queries \texttt{get(appID, h)} for the height of interest, yielding
        root $r_{\text{chain}}$.
  \item Serializes their local CRDT \texttt{DocumentSnapshot} at logical
        height $h$ and computes $r_{\text{local}} = \text{SHA-256}(\text{snapshot})$.
  \item Compares: $r_{\text{local}} = r_{\text{chain}}$ implies state
        integrity at height $h$.
  \item If $r_{\text{local}} \neq r_{\text{chain}}$, the server diverged
        from the anchored state. The user holds cryptographic evidence
        of tampering.
\end{enumerate}

\noindent
Verification is a read-only \texttt{eth\_call}---zero gas, zero latency
beyond RPC round-trip ($<50$ms to a local node, $<200$ms cross-region).

% ============================================================================
\section{Adversary Model}
\label{sec:adversary}

\begin{theorem}[Anchor Integrity]
Let $\mathcal{A}$ be an adversary controlling the Base server but not
$\geq 1/3$ of chain validators. $\mathcal{A}$ cannot produce an anchor
$(appID, h, r')$ for $r' \neq r$ at height $h$ if an honest anchor
$(appID, h, r)$ was already finalized.
\end{theorem}

\begin{proof}
The height monotonicity invariant prevents overwriting: once height $h$ is
anchored, no transaction with height $\leq h$ for the same appID is
accepted. Reverting the anchor requires a chain reorganization past the
Quasar finality depth, which requires compromising $\geq 1/3$ of
validators across all three signature schemes (BLS, ML-DSA, Ringtail).
\end{proof}

\paragraph{Withholding attack.}
The server can refuse to anchor, or anchor selectively. This is detectable:
if the user expects an anchor at height $h$ and \texttt{getLatest(appID)}
returns $h' < h$, the server is withholding. Mitigation: let clients
submit anchors directly (the precompile accepts any \texttt{msg.sender}).

\paragraph{Equivocation.}
The server anchors root $r_1$ at height $h$ on-chain but presents root
$r_2$ to a client. The client detects this on verification (step 4 above).
The on-chain record is authoritative.

\paragraph{Reorg protection.}
On Q-Chain, a QuasarCert-finalized anchor is irreversible under the Quasar
security model. On C-Chain, wait for 2 block confirmations ($\approx 4$s)
for probabilistic finality equivalent to $10^{-9}$ reorg probability.

% ============================================================================
\section{Cost Analysis}
\label{sec:cost}

\begin{center}
\begin{tabular}{lrrr}
\toprule
\textbf{Chain} & \textbf{Gas/anchor} & \textbf{Cost (LUX)} & \textbf{Cost (USD)} \\
\midrule
C-Chain & 25,400 & 635 nLUX & \$0.001 \\
Q-Chain & $\approx 0$ & $<1$ nLUX & $<$\$0.000001 \\
\bottomrule
\end{tabular}
\end{center}

\paragraph{Frequency trade-off.}
Anchoring every 5 minutes = 288 anchors/day = \$0.29/day on C-Chain,
essentially free on Q-Chain. Anchoring every 1,000 operations at 100
ops/sec = anchor every 10 seconds = 8,640/day = \$8.64/day on C-Chain.

\paragraph{Recommendation.}
Default to every 5 minutes \emph{and} every 1,000 operations (whichever
comes first). For cost-sensitive deployments, use Q-Chain. For EVM-native
apps already on C-Chain, use C-Chain to avoid cross-chain queries.

% ============================================================================
\section{Batched Anchoring}
\label{sec:batching}

For applications managing many documents (e.g., a multi-tenant Base
deployment), individual anchors per document are wasteful. Instead:

\begin{enumerate}[nosep]
  \item Collect up to 1,000 per-document Merkle roots.
  \item Build a Verkle tree (or SHA-256 Merkle tree) over these roots.
  \item Submit one anchor: $(appID_{\text{batch}}, h_{\text{batch}}, r_{\text{batch}})$.
  \item Publish the Verkle proof for each document's root as an off-chain
        inclusion witness.
\end{enumerate}

\noindent
This amortizes the 25,400 gas cost across 1,000 documents: 25.4 gas per
document. The off-chain Verkle proof is $\approx 1$\,KB and can be stored
alongside the CRDT state.

Batching is a client-side optimization. The precompile does not change;
it anchors one root per call regardless of whether that root represents
one document or a thousand.

% ============================================================================
\section{Implementation}
\label{sec:implementation}

\subsection{Precompile}

The precompile is implemented as a Go module at
\texttt{github.com/luxfi/precompile/anchor}, registered in the Lux
precompile registry. It implements
\texttt{contract.StatefulPrecompiledContract} with four ABI-dispatched
functions.

Storage keys use \texttt{keccak256(appID || height)} for root storage
and \texttt{keccak256(appID || 0xFF..FF)} for the latest-height slot.
The sentinel \texttt{0xFF..FF} avoids collision with any valid height
(uint64 max is $2^{64}-1$, which fits in 8 bytes and does not collide
with the 32-byte sentinel).

\subsection{Client SDK}

The client SDK at \texttt{github.com/hanzoai/base/crdt} provides:

\begin{lstlisting}
type Anchorer interface {
    Submit(ctx context.Context, root [32]byte) error
    Verify(ctx context.Context, height uint64, root [32]byte) (bool, error)
}
\end{lstlisting}

An \texttt{RPCAnchorer} implementation talks to the chain via JSON-RPC.
A background goroutine anchors on a configurable interval.

% ============================================================================
\section{Comparison with Prior Art}
\label{sec:comparison}

\begin{center}
\begin{tabular}{lccccc}
\toprule
 & \textbf{LP-Anchor} & \textbf{Ceramic} & \textbf{KERI} & \textbf{Radicle} & \textbf{Nostr} \\
\midrule
PQ-safe finality & Yes (Q-Chain) & No & No & No & No \\
On-chain storage & 2 slots/anchor & Tx log & Off-chain & Tx log & None \\
Batching & Verkle tree & Merkle tree & KEL & Git tree & N/A \\
Verify latency & $<200$ms & Seconds & Seconds & Seconds & N/A \\
Gas/anchor & 25.4k & $\sim$50k & N/A & $\sim$30k & N/A \\
\bottomrule
\end{tabular}
\end{center}

\noindent
Ceramic anchors CIDs to Ethereum via a Merkle tree of stream tips,
batching every 30 minutes. KERI maintains key event receipt logs off-chain
with witness infrastructure. Radicle and Gitopia anchor Git commit hashes.
Nostr event IDs are self-certifying but have no chain finality.

LP-Anchor is the only design with post-quantum finality guarantees and
sub-second verification latency.

% ============================================================================
\section{Conclusion}
\label{sec:conclusion}

LP-Anchor provides a minimal, high-assurance mechanism for anchoring
CRDT state to a public blockchain. The design requires 72 bytes per anchor,
25,400 gas on C-Chain (near-zero on Q-Chain), and delivers sub-200ms
verification latency. The height monotonicity invariant and Quasar finality
guarantee that anchored state cannot be retroactively altered without
compromising the chain's validator set across three independent
cryptographic assumptions.

The precompile is general: any application that computes a periodic state
hash can use it. Base CRDT anchoring and \texttt{liquidity/state}
reconciliation anchoring are the first two consumers.

\bibliographystyle{plain}
\begin{thebibliography}{9}

\bibitem{quasar}
Z.~Kelling, ``LP-105: Quasar --- Post-Quantum BFT Consensus,''
\textit{Lux Research}, April 2026.

\bibitem{crdt}
M.~Shapiro et al., ``Conflict-free Replicated Data Types,''
\textit{Proc. SSS}, 2011.

\bibitem{ceramic}
Ceramic Network, ``Ceramic Protocol Specification v2,''
\url{https://ceramic.network}, 2024.

\bibitem{keri}
S.~Smith, ``Key Event Receipt Infrastructure (KERI),''
\textit{IETF Draft}, 2022.

\bibitem{eip2929}
V.~Buterin, M.~Swende, ``EIP-2929: Gas cost increases for state access opcodes,''
\textit{Ethereum Improvement Proposals}, 2021.

\end{thebibliography}

\end{document}
